\section{Results}

% claims: C08; sources: S10 S11
The primary end-to-end result is that all 128 public closures passed the frozen
integrated endpoint. In every matched pair, the output-side signal followed the
authority derived from the authenticated state while the prompt, answer,
surface payload, configuration, output indices, and primitive random draws
remained paired.

\paragraph{Causal and receipt evidence.}
% claims: C03 C04 C05; sources: S04 S05 S06
Before integration, the selected model passed its prospective causal
confirmation. Natural state and answer predictions were correct on 512/512
oriented examples, and every registered answer-preserving and
answer-changing intervention family achieved 512/512 complete joint success.
The learned transition continued to generate downstream \zthree{} after
\ztwo{} interventions. A separate receipt confirmation verified the exact
event, abstract state class, and registered chain distinctions under its frozen
protocol tests. These results concern the selected model and trusted verifier,
not a naturally discovered representation or trust-free attestation.

\paragraph{Generic carrier.}
% claims: C07; sources: S09
The prospective generic-carrier validation passed 512/512 closures at the
frozen operating point. There were no wrong-authority familywise events and one
carrier-off familywise event within its registered gate. This establishes
reliable carriage for the tested 16-authority family and grammar, not a robust
four-bit capacity result or receipt binding by itself.

\begin{table}[t]
% claims: C08 C10 C11; sources: S10 S14 S15 S16 S17
\caption{Primary prospective outcomes. A closure is the independent matched
unit; nested outputs and checks are not replications. The two protected attack
failures were consequences of the single computation failure.}
\label{tab:results}
\centering
\small
\begin{tabularx}{\linewidth}{@{}YccY@{}}
\toprule
Role & Primary unit & Outcome & Interpretation \\
\midrule
Public integrated validation & closure & 128/128 & all frozen gates passed \\
Protected qualification & closure & 127/128 descriptive & not qualified; one computation failure \\
Protected confirmation & closure & not run & sealed and not tested \\
\bottomrule
\end{tabularx}
\end{table}


\paragraph{Public integration.}
% claims: C08; sources: S10 S11
The public receipt-bound integration was prospectively validated on 128/128
held-out public closures, with zero registered attack failures, zero
carrier-off regressions, and zero wrong-authority confusions in that
population. Across the 256 natural/alias execution packages nested within those
closures, authenticated aggregate scores had median 10.54 and range
$[8.66,12.27]$ against the frozen threshold $T=4.04$; the
authenticated-minus-strongest-wrong margin had median 8.88 and range
$[6.35,10.83]$. These score summaries are post-hoc descriptions of frozen rows,
not additional prospective endpoints or independent replications.

% claims: C08; sources: S10 S11
A post-hoc 95\% Wilson interval for the 128/128 closure count is
$[0.9709,1.0000]$ if closures are treated as exchangeable Bernoulli units. It
does not establish population representativeness, deployment reliability, or
a universal false-positive rate.

\paragraph{Protected non-qualification.}
% claims: C10; sources: S14 S15
Protected qualification was not earned: one of 128 closures failed computation
fidelity, the descriptive result was 127/128, and final protected confirmation
was not run. The same upstream failure caused two dependent attack-expectation
failures. Because the computation and zero-attack gates were
non-compensatory, the protected population did not qualify despite the other
127 closures. The cause cannot be diagnosed or repaired from the closed
evidence.

% claims: C11; sources: S16 S17 S20
The separately registered protected confirmation population contains 256
sealed closures but was never authorised or executed. It therefore provides no
receipts, outputs, detector decisions, or outcome evidence.

\paragraph{Preserved negatives.}
% claims: C12; sources: S18 S21
The frozen training recipe passed only two of three seeds; Model 2 failed a
registered natural-$\ztwo=0$ subgroup gate. An earlier carrier operating point
also failed its advancement rule at 88/128 pairs. These negatives motivate the
selected-model claim and the final carrier design; neither is pooled with the
integrated result.
